\chapter{模型评测}
\label{chap:nlp-evaluation}

一套语言模型在固定试题上答对得更多，可能来自参数学习，也可能来自更长的提示、更多次
采样或一个更强的选择器。若把这些改变同时发生的结果统称为“模型能力提升”，我们就
无法判断下一阶段该改数据、训练还是推断。本章承接前面各阶段，研究怎样把观察转成
有范围的证据：固定被测对象，准备独立问题，规定判定办法，再报告差异及其不确定性。

评测（evaluation）是按预先规定的材料与规则观察系统行为；它不等于给一个总分。
我们沿用通知、数学求解与多轮交流的例子，同时区分基础概率模型、交互模型和带检索或
工具的完整应用。前几章的阶段检查用于决定下一轮训练，本章给出这些检查共同依赖的
协议；后续应用章再展开跨度、译文、证据记录和执行状态的具体标准。

\section{评测对象与协议设定}
\label{sec:nlp-eval-protocol}

\subsection{模型、推断过程与应用系统}
\label{sec:nlp-eval-unit}

把同一组参数接到两种解码器上，回答准确率可能改变。因此，被测对象至少应说明参数版本、
词元接口、提示与对话模板、解码和停止规则。检索器、外部资料、工具、验证器参与回答时，
它们也属于产生结果的过程。只给模型名称不足以恢复一次比较。

可以把协议记作$\mathcal P=(\bm\theta,\mathcal D,\mathcal I,\mathcal B,\mathcal J)$：
$\bm\theta$是模型参数，$\mathcal D$是被测请求的分布或样本，$\mathcal I$是输入与推断
过程，$\mathcal B$是预算约束，$\mathcal J$是判定器。这个记号用来固定比较条件，并不
假定五者相互独立。研究训练改变时尽量固定后四项；研究推断改变时固定参数并按相同预算
比较。若评测完整应用，就报告组合结果，同时用组件替换定位收益来源。

\subsection{目标行为与可观察判据}
\label{sec:nlp-eval-criteria}

对通知提问“校外来宾何时截止报名”，正确日期、时区和对象范围可由材料直接检查。
若要求“写得亲切一些”，则存在多个合法表达，单一参考字符串不能覆盖全部答案。
我们需要把要求拆成可判定的条件：日期是否正确、是否遗漏截止时间、是否添加未经支持的
承诺，以及语言是否适合预定读者。条件可以分别报告，也可以预先规定怎样合并。

设一条请求有$K$个必要条件，条件判定为$j_k(x,y)\in\{0,1\}$。单条请求的严格合格指示为
$\prod_{k=1}^K j_k$，因为任何必要条件失败都会使该请求不合格；在题集上对这个指示取平均，
才得到严格合格率。单条请求的逐条件完成比例为$K^{-1}\sum_kj_k$，
再对题集平均，反映通常完成了多少条件。两种聚合不能互换。假设每个回答都遗漏同一项
重要截止条件，逐条件平均可能很高，严格合格率却为零。这不是指标矛盾，而是验收对象
不同。HELM以多场景、多指标报告语言模型行为，为避免单分数掩盖差异提供了具体参照
\scite{nlp-intro-liang2023helm}。

HELM把一个场景表示为任务、领域和语言的组合，适配器把不同模型接口转成同一评测请求，再由各项指标读取原始回答。其七类核心指标包括准确性、校准、稳健性、公平性、偏差、毒性和效率；这些指标不天然等权，也不都能用于每个场景。场景、提示和聚合是协议设计，不需要训练一个“HELM模型”；若某项指标调用毒性分类器，则应另外说明该分类器的参数及标注来源。

原报告在30个模型、16个核心场景和26个专项场景上开展评测，将模型在核心场景中的平均可比覆盖从此前报告的$17.9\%$提高到$96.0\%$；112个核心场景与指标组合中实际测量98个。这里改善的是评测覆盖，不是模型正确率。报告还观察到，同一模型族中扩大参数常提高准确率并增加成本，但不同族之间不存在一个统一的准确率与效率排序。HELM的实际用途是为模型选择提供同协议、多指标证据，其初始核心范围主要是英语，不能直接代表中文、多轮工具执行或所有安全行为。

\subsection{预算、版本与复现记录}
\label{sec:nlp-eval-budget}

推理专题已经表明，多采样能增加找到正确候选的机会。比较时应同时记录输入长度、允许
生成长度、候选数、验证次数、工具调用和实际消耗。相同词元上限不一定意味着相同时间或
费用，不同词元化下的“同样一千词元”也未必包含相同文字。报告应说明选择哪种预算，
以及为何它对应当前使用条件。

除了平均用量，还应记录超限、无结束符、解析失败和拒绝回答的处理规则。只对成功返回的
样本统计时延会遗漏超时；把格式失败悄悄重试到成功，会增加未披露的预算。保留模型与数据
版本、提示、随机种子和逐请求记录，可以重放具体失败。固定随机种子有助于复现，但单次
随机运行仍不足以说明结果对种子的稳定程度。

\section{评测材料与独立性}
\label{sec:nlp-eval-data}

\subsection{划分单位与污染检查}
\label{sec:nlp-eval-contamination}

预训练、SFT、偏好反馈和提示开发都可能接触评测题。即使测试文件没有直接加入训练集，
同一篇通知的改写、答案解析和另一语言版本也可能已经出现。因此，划分应按源文档、事件、
模板或问题族执行，再检查精确重复与语义近邻。逐行随机划分往往不足以阻断共同来源。

污染检查没有发现重复，只能说明指定规则未找到重叠；不能证明模型从未接触过题目。
可辅以训练时点之后的新材料、保密题、动态生成且人工核验的问题，以及结构相同但内容
改变的对照。新题也要保证难度和标注质量可比，否则分数下降既可能来自少记忆，也可能
来自题本身更难。数据独立性与数据可比性需要同时检查。

\subsection{能力覆盖与条件分组}
\label{sec:nlp-eval-coverage}

一套中文通知问答不能代表数学、代码或其他语言的能力；一套选择题也不能自动代表多轮
办事。应从拟支持的请求范围出发，按语言、领域、长度、难度、信息可见性和交互轮次
分组，再说明每组样本如何取得。公开基准用于可重复比较，场景样本用于检查目标使用
条件，两者可以互补，却不能仅靠数量相加就消除覆盖缺口。

分组还影响总分。设组$g$的准确率为$\widehat a_g$、权重为$\omega_g\geq0$且$\sum_g\omega_g=1$，总分为
$\sum_g\omega_g\widehat a_g$。按测试样本量加权反映当前题集组成，按目标请求频率加权
才对应指定使用分布；对各组等权又是另一种问题。权重改变会使模型排序改变，应同时报告
组内结果及权重依据，而不是只保留一个平均数。

\subsection{对照题与行为探针}
\label{sec:nlp-eval-probes}

为了判断模型是否使用了通知中的限制条件，可以只把“校外来宾”改为“全体参加者”，
观察答案是否相应改变。也可以保持句意而替换地名、日期格式或语序，检查不该改变的
判断是否保持。这类成对题控制了一个因素，比仅看题集总分更容易定位失败。

对照成立的前提是改动没有暗中改变其他验收条件。把时间移到过去可能同时改变可执行性，
把人物换成不同身份可能改变授权。行为探针说明模型在特定干预下怎样响应，不直接揭示
内部因果机制。若把输出变化解释为某个注意头或训练阶段的作用，还需要相应干预证据。

\section{概率、任务结果与偏好判定}
\label{sec:nlp-eval-measures}

\subsection{留出似然与概率质量}
\label{sec:nlp-eval-likelihood}

预训练阶段最直接的观察是模型给未参与更新的词元分配多少概率。固定非空的计分位置集合
$\mathcal M$，以自然对数计算平均负对数似然与\term{困惑度}：
\begin{equation}
 H=-\frac{1}{|\mathcal M|}\sum_{(i,t)\in\mathcal M}
 \log p_{\bm\theta}(w_{it}\mid c_{it}),\qquad \operatorname{PPL}=\exp(H).
 \label{eq:nlp-eval-ppl}
\end{equation}
$c_{it}$是该位置实际可见的条件。数值越小，参考词元概率的几何平均$\exp(-H)$越高；这不等于算术平均概率更高。
直接比较时，词元化、上下文截断、可见条件、掩码与计分单位须保持相同约定。
模型结构可以不同，只要计算的是相同条件下同一参考事件的概率。

词元化不同会改变分母与预测事件。跨词元接口时，可以在同一原始文本、同一覆盖范围和
明确边界条件下报告每字节负对数似然，但仍需说明字节编码及哪些开销计入。掩码恢复模型
的某个去噪损失也不自动等于自回归似然，若使用似然界或估计器，应单列其定义。
概率质量可以帮助筛选预训练参数，却不能证明模型遵守指令、答案正确或使用证据可靠。

计算留出似然时，模型参数冻结，参考序列只提供计分目标及合法的先前条件，不反向更新模型。长文本可滑动窗口读取已有前缀，重叠区域用于条件但不重复计分；在互不相连的小块中逐块重启上下文，则测量的是另一个条件概率。对选择题计算选项似然，也须声明是否含选项标签、是否按长度归一化，以及同义选项怎样处理。长度归一化改变排序目标，不能把“平均每词元概率高”自动写成“完整选项概率高”。

双向掩码模型没有普通左到右联合似然时，可以逐位置遮蔽参考词元，再用其余位置预测该词元，形成伪对数似然
\begin{equation}
 \operatorname{PLL}(w_{1:T})
 =\sum_{t=1}^{T}\log p_{\bm\theta}(w_t\mid w_{1:t-1},w_{t+1:T}).
 \label{eq:nlp-eval-pseudo-likelihood}
\end{equation}
每一项的条件不同，不能把这些条件概率的乘积默认当作归一化序列分布。PLL不新增训练参数，但通常需要多次掩码前向；若另训一个单次评分器去逼近PLL，就必须报告新的训练目标。Salazar等在固定LibriSpeech的100条转写候选上，以开发集选择语言评分权重，测试集重排；表2（第2703页）中RoBERTa-large把test-clean的词错误率从$7.26\%$降至$5.05\%$，test-other从$20.37\%$降至$16.79\%$。这展示了语言可接受性评分的实际用途，不能解释为语音模型参数提高，也不能据此把PLL和自回归PPL混为同一指标\scite{nlp-eval-salazar2020masked}。

\subsection{答案核验与生成质量}
\label{sec:nlp-eval-answer}

数学题有可规范化的数值答案，代码有可执行测试，通知问答有来源字段。判定器先约定答案
提取与规范化规则，再计算正确率；提取失败应有独立计数，不能在看过模型输出后逐条放宽
规则。单元测试通过说明覆盖的输入输出条件成立，未覆盖路径和规格仍可能出错。

开放生成需要同时检查信息与表达。字符串重合、语义相似、事实蕴含和人工可读性评判测量
的对象不同：更相似的句子可能保留同一个错误日期，更自由的合法改写也可能得不到高重合
分。应把关键事实与硬约束作为独立判据，再用替代指标观察表达差异。各类生成指标的
机制及局限在文本生成与转换章展开，本章只规定不能把代理分数直接等同于任务合格。

对可验证约束，可以在每条提示中列出多项要求，并分别核验词数、格式、关键词等条件。
IFEval是这一思路的代表：以可检查的指令约束减少完全依赖开放式主观评分的情形
\scite{nlp-eval-ifeval2023}。格式判定通过仍不保证语义正确，验证器自身也应有单独测试。

IFEval把提示拆成约束身份与参数，例如“回答恰含三段”或“不得出现指定词”，模型生成完整回答后，由确定程序逐项判定。原版本包含25类可验证指令与541条提示，一条提示可以组合多项约束。评测器不拟合语言模型参数，关键词、段数及语言识别等检查器却须有明确的解析约定；中文的“词数”不能未经说明套用英语空格分词。提示级严格指标要求所有约束通过，指令级严格指标只按每项统计，对应本章前述两种验收对象。

宽松指标先移除常见Markdown修饰或首尾附加语，再检查变换后的文本；它可减少格式误判，也可能使原本违反词数的回答变成合格，不能替代严格值。原论文表3（第5页）中GPT-4的提示级严格合格率为$76.89\%$、指令级为$83.57\%$，宽松口径分别为$79.30\%$和$85.37\%$。这组同模型结果说明验收单位和文本变换本身就会改变分数。IFEval适用于输出约束保持和格式回归检查，不检查通知日期是否有来源，也不保证所有约束类型都符合真实使用的重要性。

\subsection{人工比较与模型评判}
\label{sec:nlp-eval-judges}

偏好对齐追求的是在指定判断标准下更常被选择，而不是提高奖励器自己的分数。独立评价
应另取提示、回答和评判者，隐藏模型身份，并平衡候选展示次序。若$A$胜、$B$胜和平局
数量分别为$W,L,T$，一种预先声明的胜率为$(W+T/2)/(W+L+T)$；同时报告三类计数，
使平局处理可审计。换一套判断标准或评判人群，就可能得到不同排序。

用大模型评判（LLM-as-a-judge）可以降低逐条人工阅读的成本，但评判模型也是一个带偏差
的预测器。MT-Bench与Chatbot Arena的研究考察了其与人工偏好的关系，也讨论了位置、
冗长及自我偏好等问题\scite{nlp-intro-zheng2023judge}。因此，应以人工标注子集校准
当前任务上的一致性，交换候选顺序检查位置影响，并记录评判提示及版本。曾参与训练的
奖励器可以用于诊断，不能单独充当独立改进证据。

模型评判可以采用单回答评分，也可以把同题的两条回答并列比较。前者需要稳定的评分尺度，后者需要固定平局、次序与比较标准；多轮题还须把先前问答交给判定器，否则它无法检查本轮是否保持前文要求。现成GPT-4评判的参数来自其已有训练，本次只是条件生成分数或胜负；人工校准子集用于验证或选择提示，不能因此宣称已经训练了独立奖励模型。若用这些标签再微调评分器，训练题、损失及测试划分须另行说明。

MT-Bench采用80个多轮问题覆盖八类任务，原研究用专家判断和Arena的用户偏好检查模型评分。表5（第8页）在排除平局的S2口径下，GPT-4成对评判与人的一致率约$85\%$，包含平局等结果的S1口径约$66\%$；前者不能简称所有回答上的85%正确。表2（第5页）交换近似回答的位置时，默认GPT-4只有$65.0\%$保持一致，说明总体一致性高也不能免除位置偏差。它适合开放表达的批量比较和初步筛查，事实、算术或动作授权仍应使用相应证据。

评判能力还可能低于独立解题能力：原研究发现GPT-4可单独答对某些算术题，却被候选中的错误说明带偏；参考答案辅助判定能减少这类错误，但又改变了判定器可见信息。关于自我偏好，论文只观察到部分判定器相对人工更偏向某些模型，未在有限数据下确定这就是因果意义的自我增强。不能将这种受限观察写成所有模型必然偏爱自身。交换次序、抽查事实和保留分歧，分别处理不同失败，不应只以换一条评判提示作为通用修复。

表~\ref{tab:nlp-eval-observation-contracts}把几种评测方法放回同一请求与输出集合。它们没有一个共同的“最高分最好”尺度；相同候选预算也不能让预测概率、约束合格与偏好票成为同一事件。

\begin{table*}[htbp]
 \begin{tabularx}{\linewidth}{@{}p{25mm}p{34mm}p{39mm}X@{}}
  \toprule
  方法 & 实际输入 & 参数与判定来源 & 验收对象及不能替代的证据\\
  \midrule
  留出似然 & 参考词元与可见条件 & 冻结语言模型；计分无训练 & 参考概率；不验指令或事实\\
  PLL重排 & 固定候选、逐位置掩码 & 冻结掩码模型；权重开发集选 & 语言可接受性；非联合似然\\
  HELM & 统一场景与模型适配 & 预设协议；部分指标用分类器 & 多指标画像；不代表全部语言与任务\\
  IFEval & 约束及完整回答 & 确定检查规则，无模型更新 & 逐约束或整提示合格；不验内容真实性\\
  人工比较 & 盲化同题候选 & 标注者及判断标准 & 目标人群偏好；仍可分歧\\
  模型评判 & 问题、回答及评分提示 & 现成模型或另训评分器 & 代理偏好；必须校准误判与偏差\\
  \bottomrule
 \end{tabularx}
 \caption{概率、规则、综合协议与偏好评价的观察对象。固定任务、输入、候选预算和判定信息后，才能解释每项证据支持的结论。}
 \label{tab:nlp-eval-observation-contracts}
\end{table*}

\section{推理、稳健性与安全行为}
\label{sec:nlp-eval-reasoning-safety}

\subsection{答案能力与轨迹可靠性}
\label{sec:nlp-eval-reasoning}

推理轨迹长度增加后，最终答案可能更常正确，也可能只是更常尝试。一次生成的正确率、
多候选中至少存在一个正确答案的比例，以及选择器最终交付的正确率，应分开统计。
第一个观察单次生成，第二个观察候选覆盖，第三个还包含验证器的能力。

设一道题独立同分布采样$n$次，其中$c$次通过固定验证器。对$n\geq k$，常用的
pass@$k$估计为
\begin{equation}
 \widehat{\operatorname{pass@}k}
 =1-\frac{\binom{n-c}{k}}{\binom nk},
 \label{eq:nlp-eval-passk}
\end{equation}
其中$n-c<k$时分子为零。式子计算从$n$条记录不放回抽$k$条时至少有一条通过的概率；
在独立同分布采样假设下，它可估计$k$次生成的候选覆盖\scite{nlp-intro-chen2021}。
例如$n=5,c=2,k=2$时为$1-3/10=0.7$。它没有规定怎样从两条候选中选对答案，
不能作为实际选择器的交付准确率。相关采样或自适应搜索还会改变其统计解释。

答案正确也不足以证明每一步解释可靠。应另检查中间结论、引用、程序执行记录与反事实
改题后的响应；从外部观察到的书面轨迹，不应直接当作全部内部计算的忠实记录。
报告推理能力时，答案、过程可检验性和成本分别提供证据，缺一项就限制相应结论。

\subsection{有害响应与过度拒绝}
\label{sec:nlp-eval-refusal}

安全对齐同时面临两种失败：对明确禁止的请求提供有效帮助，对合法请求作出不必要拒绝。
分别准备按既定政策标注的风险请求集$\mathcal H$和正常请求集$\mathcal S$，报告有害帮助
比例和正常帮助成功率。若只统计风险集上的拒绝比例，一个对所有请求都拒绝的模型也会
显得优秀；若只统计正常帮助，又会漏掉风险边界。

XSTest通过安全请求及容易触发误判的表述检查过度安全行为，提示我们拒绝本身也需要
与请求含义对应\scite{nlp-eval-xstest2024}。对于允许解释但不允许执行的请求，还应区分
回答边界与动作边界。评判要依据明确政策、上下文和用途，关键词命中不能单独决定风险。
政策改变后应更新判定规则并保留版本，不能把两个政策下的拒绝率直接解释为模型进退。

XSTest并非训练一个更安全的模型，而是构造含有风险词却语义合法的请求，并配对真正不安全的请求来检查边界。其公开套件含250条安全提示、200条不安全对照；安全提示包括比喻、游戏语境、公开人物信息等类型。人工把回答分为完整拒绝、部分拒绝与遵从，再按请求类型汇总，拒绝检测器只是辅助标注，不能把所有拒绝语句出现都算完整拒绝。

原论文表1（第5379页）在250条安全提示上，Llama-2-70B-Chat原系统提示下完整拒绝为$38.0\%$、部分拒绝为$21.6\%$；Mistral-7B-Instruct的两类拒绝各为$0.8\%$。同一Mistral接入安全系统提示后，完整和部分拒绝又分别为$9.6\%$和$9.2\%$。这些观察既展示提示属于被测对象，也说明低安全集拒绝率不能单独证明风险控制更好；还须看配对不安全集的行为。实际用途是诊断过度拒绝和安全关键词捷径，不是用小套件估计全部真实风险的发生概率。

\subsection{对抗条件与输入边界}
\label{sec:nlp-eval-adversarial}

正常提问通过以后，还要问：攻击者能改变哪些输入，使模型偏离已有要求？越狱评测通常
研究攻击提示诱导被禁止响应；提示注入则常把指令放在检索文档、网页或工具返回中，试图
让模型把低信任材料当作有权修改任务的命令。两者的可控位置、可见信息和成功条件应
分别说明。仅在用户提示上测试，不能覆盖外部材料引起的越权。

一次攻击评测要固定攻击者可访问的信息、尝试次数、反馈及判定器。JailbreakBench用
标准化威胁模型、提示与评分规则支持可比较的越狱实验\scite{nlp-eval-jailbreak2024}。
攻击失败率只对声明的攻击集合和预算成立；更强攻击或新接口仍需重新检查。稳健性与
安全性的通用理论可参考\BookChapterRef{chap:trust-security}{对抗安全}，本章把要求落实
到语言输入、上下文边界和可观察输出上。

JailbreakBench固定目标模型的系统提示、聊天模板、攻击者可改的用户输入及查询预算，攻击算法提出提示、读取允许的反馈，再由判定器检查回答是否提供了所要求的有害帮助。原始JBB-Behaviors包括十类共100种有害行为；按行为统计至少一次攻击成功，与按所有尝试统计成功率有不同分母。搜索器和攻击字符串可以不学习参数，梯度攻击、训练攻击模型和调用现成模型生成攻击也有不同资源，不能都记成一次普通请求。

判定器本身是重要误差来源。原论文表2（第9页）在100组提示回答、三位专家多数标签的对照中，GPT-4判定一致率为$88\%$，误报率$16\%$、漏报率$7\%$；Llama Guard一致率为$76\%$，误报率$7\%$、漏报率$47\%$。作者按控制误报选择后者，并不表示它所有指标最高，也不表示拒绝词检测等同专家标签。这套协议可用于比较固定威胁模型下的攻击和防御，但判定误差与预算限制须进入结论，不能把没判出成功写成系统没有漏洞。

提示注入的验收还要追踪信息权限。对同一通知应用，可以只替换低信任材料中的一段文字，保留用户目标与工具权限，观察是否发生未经授权的要求变更；攻击成功对象可能是泄露、错误引用或越权动作，不只是输出了风险词。低信任内容参与模型条件计算，不代表它有权改写系统任务。对抗搜索可使用固定策略或另训攻击模型，但一个仅测试用户提示越狱的实验，不能支撑检索材料和工具返回均安全的结论。

\section{统计比较与阶段诊断}
\label{sec:nlp-eval-analysis}

\subsection{差异估计与不确定性}
\label{sec:nlp-eval-uncertainty}

两个模型在同一批$n$条请求上比较时，先为每条计算差值$d_i=s_A(x_i)-s_B(x_i)$，再求
$\widehat\Delta=n^{-1}\sum_i d_i$。成对差值保留了同一请求的难度关联，比把两组分数
当作不相关样本更符合实验设计。若请求按事件或文档成组，应以这些组为抽样单位，不能
把高度相关的多个问句当成独立重复。

可以按原抽样单位重采样，报告成对bootstrap区间；其解释依赖测试样本能代表目标总体，
也依赖聚类结构处理正确。仅有固定题集时，区间反映该抽样模型下的不确定性，不覆盖未来
分布变化、未计入的提示选择或未观测风险。多次比较后只挑最高分，也会引入选择偏差，
最终确认应使用独立材料。

\subsection{阶段贡献与失败归因}
\label{sec:nlp-eval-attribution}

比较预训练起点、SFT后、偏好对齐后与推理强化后的参数时，应使用共同可运行的协议，
并另报告各阶段推荐的输入模板。共同协议帮助定位参数变化，推荐模板观察实际用法，
二者不能相互替代。基座不遵守对话模板不必然表示其语言概率建模更差，聊天模型更受偏好
也不保证领域知识增加。

要判断某项数据或算法的作用，需要只改变该因素，或清楚声明哪些因素无法分开。
多个改动一起产生的收益只能归于组合。再按错误类型追踪：缺知识时考虑继续预训练，
格式与任务响应不稳定时考虑示范，评价标准不符时检查偏好反馈，候选覆盖足而选择错误时
检查验证器。这些是由证据支持的诊断方向，最终还需干预实验检验，不能由错误名称直接
推出唯一训练方案。

\subsection{完整报告与使用范围}
\label{sec:nlp-eval-report}

一份可使用的报告应让读者恢复被测对象、题集来源、信息边界、预算、判定、组内结果与
失败记录。样本数、分母、无效响应和区间与总分一起给出，才能解释“提高了多少”。
对重要约束还需列出绝对水平：假设错误率从$20\%$降到$10\%$，下降一半不表示剩余错误
已经满足一个要求错误率低于$1\%$的用途。这是构造算例，说明相对改进与可接受阈值必须
分别回答。

报告的最后应说明哪些语言、长度、时点、任务和攻击条件已有证据，哪些尚未覆盖。
模型进入新场景后继续收集失败，但不能把新近用于调参的记录又当成独立测试。评测由此
成为训练与应用之间反复更新的证据接口，而不是模型完成后永久有效的一张成绩单。

\section{本章小结}
\label{sec:nlp-eval-summary}

语言模型的评测结论同时依赖参数、材料、推断过程、预算和判定规则。留出似然观察概率
建模，任务核验观察交付结果，偏好比较观察指定标准下的选择，安全与对抗测试观察行为
边界。各类证据可以互补，但不能用其中一项替代其余对象。

前面各阶段现在具有了共同的比较依据：改变学习环节时固定可比较条件，改变推断方法时
把额外计算计入预算，再以分组结果、反例和不确定性限制结论。后续应用章将把同一原则
落实到文本结构、知识记录、生成内容、交互行动与语音时间条件，说明一般模型怎样才算
完成一个具体任务。

\paragraph{延展阅读}
HELM适合继续研究多场景、多指标的统一报告，IFEval把部分指令要求转成可检查规则，
XSTest与JailbreakBench分别考察过度拒绝和越狱条件。对照阅读时，应先写清每项工作
评价的行为对象、攻击或请求分布、判定器和统计单位，再判断它能否支持当前用途；
单一总分不能替代能力、稳健性与安全边界的分项证据
\scite{nlp-intro-liang2023helm,nlp-eval-ifeval2023,nlp-eval-xstest2024,nlp-eval-jailbreak2024}
